\documentclass[10pt,a4paper]{amsart}
\usepackage[margin=19mm]{geometry}
\usepackage{amsmath,amssymb,mathtools}
\usepackage[scr=rsfs,cal=euler]{mathalfa}
\usepackage{xfrac}
\usepackage[unicode,hidelinks,bookmarks=false]{hyperref}
\newcommand{\ie}{i.e.\ }
\newcommand{\cf}{cf.\ }
\newcommand{\resp}{respectively\ }
\newcommand{\Subsection}{\subsection}
\providecommand{\nobreakdash}{\nobreak\hbox{-}}
\setlength{\emergencystretch}{2em}
\multlinegap0pt
\usepackage{bm}
\usepackage{bbm}
\usepackage{dsfont}
\usepackage{xspace}
\usepackage{cancel}
\usepackage{mathtools}
\usepackage{amsthm}
\makeatletter
\@ifundefined{theorem}{\newtheorem{theorem}{Theorem}}{}
\@ifundefined{lemma}{\newtheorem{lemma}[theorem]{Lemma}}{}
\makeatother
\renewcommand{\le}{\leqslant}
\title[Every expansive $ m $-concave operator has $ m $-isometric d]{Every expansive $ m $-concave operator has $ m $-isometric dilation}
\author{Micha{\l} Bucha{\l}a}
\date{Source: arXiv:2507.07252v1 (CC BY 4.0)}
\begin{document}
\maketitle

\noindent\emph{Source and licence.} Every expansive $ m $-concave operator has $ m $-isometric dilation. Version arXiv:2507.07252v1 (CC BY 4.0), distributed under the Creative Commons Attribution 4.0 International licence, as recorded in the arXiv licence metadata for this exact version and shown on the abstract page https://arxiv.org/abs/2507.07252v1. Full rights notice: licenses/arxiv-2507.07252-CC-BY-4.0.txt.\par\smallskip

\noindent\textit{Editorial scope.} Counted unit: the Theorem whose source label is \texttt{ThmDilationOf3ConcaveWithoutExpansivity} in the article (source0.tex lines 333--336, source label \texttt{ThmDilationOf3ConcaveWithoutExpansivity}), delivered together with the article's own complete original proof of it (source0.tex lines 337--362, one \texttt{proof} environment of 26 lines). No mathematics is written, repaired or re-derived: statement and proof are the article's own text line for line. Required context retained uncounted: ThmDilationOfMConcave (lines 163--178); FormDilationMatrix (lines 166--176); LemMIsometricityOfDilation (lines 194--206); FormDilationMatrixMIsometric (lines 200--204). Every definition, display and label of those lines is retained unchanged. Omitted from this package and not redistributed: 1--162, 177--193, 205--332, 363--428. Nothing inside a retained excerpt is omitted or altered: every retained body is delivered exactly as the article's lines are written, so no omission receipt of any kind accompanies this package. Citations: the retained excerpts carry no citation command at all, so no bibliography entry of the article is transcribed and no reference is redistributed. Reference adaptation: none was needed and none was made; every reference inside the delivered unit resolves inside it, so no unresolved reference or citation is produced. Printed slips kept literal: no printed word, symbol, hypothesis or number of the counted statement or of its proof was altered. Layout and class: the article's own class is \texttt{amsart} with packages including inputenc, amssymb, babel, mathrsfs, euscript, orcidlink, cite, amsaddr, amsthm, amsfonts, amsbsy, enumitem, xcolor; the wrapper is the lane's pinned \texttt{amsart} editorial wrapper at 10pt on A4, which restates the theorem-like environments the retained text needs and adds the article's own macro definitions that the retained text needs (\texttt{\textbackslash le}), with extra wrapper packages (bm, bbm, dsfont, xspace, cancel, mathtools). The delivered numbering is the unit's own. Assets: no retained span contains a figure environment, a graphics-inclusion command, a PDF-page inclusion, a table or a TikZ picture; no image file is copied into this package, the package declares no asset dependency and no image file is redistributed with it. Licence: version arXiv:2507.07252v1 of this article is distributed under the Creative Commons Attribution 4.0 International licence, as recorded in the arXiv licence metadata for this exact version and shown on the abstract page https://arxiv.org/abs/2507.07252v1. A full rights notice is delivered at licenses/arxiv-2507.07252-CC-BY-4.0.txt. Characters: every retained excerpt is pure ASCII, so the delivered engine, XeLaTeX with its default Latin Modern text font, reports no missing character.
\begin{theorem}
    \label{ThmDilationOfMConcave}
    Let $ H $ be a Hilbert space and let $ T\in \mathbf{B}(H) $ be expansive and $ m $-concave ($ m\in \mathbb{N}_{2} $). Then $ T $ has an $ m $-isometric dilation $ W $ on the space $ H\oplus \ell^{2}(\mathbb{N}, H^{\prime}) $, where $ H' \subset H $ is a closed subspace of $ H $, with the following matrix representation
    \begin{equation}
        \label{FormDilationMatrix}
        W = \begin{bmatrix}
            T & 0 & 0 & 0 &  \ddots\\
            U & 0 & 0 & 0 &  \ddots\\
            0 & S_{1} & 0 & 0 &  \ddots \\
            0 & 0 & S_{2} & 0 &  \ddots\\
            0 & 0 & 0 & S_{3} &  \ddots\\
            \ddots & \ddots & \ddots & \ddots & \ddots
        \end{bmatrix},
    \end{equation}
    where $ U\in \mathbf{B}(H,H^{\prime}) $ is nonnegative and $ S_{j}\in \mathbf{B}(H^{\prime}) $ is positive and invertible for every $ j\in \mathbb{N}_{1} $. Moreover, $ U $ can be chosen to satisfy the equality $ T^{\ast}U^{2}T = U^{2} $ and $ S_{1},\ldots,S_{m-2} $ can be chosen to be identity operators on $ H^{\prime} $.
\end{theorem}

\begin{lemma}
    \label{LemMIsometricityOfDilation}
    Let $ W\in \mathbf{B}(H\oplus\ell^{2}(\mathbb{N},H^{\prime})) $ be as in \eqref{FormDilationMatrix}. For $ m \in \mathbb{N}_{1} $ the following conditions are equivalent:
    \begin{enumerate}
        \item $ W $ is $ m $-isometric,
        \item the unilateral weighted shifts $ S\in \mathbf{B}(\ell^{2}(\mathbb{N},H^{\prime})) $ with weights $ (S_{j})_{j\in \mathbb{N}_{1}} $ is $ m $-isometric and
        \begin{equation}
            \label{FormDilationMatrixMIsometric}
            %\sum_{k = 0}^{m}(-1)^{m-k}\binom{m}{k} \lVert T^{k}h\rVert^{2}
            \langle \beta_{m}(T)h,h\rangle +\sum_{\ell=1}^{m}(-1)^{m-\ell}\binom{m}{\ell}\sum_{k=1}^{\ell} \lVert S_{k-1}\cdots S_{1}UT^{\ell-k}h\rVert^{2} = 0.
        \end{equation}
    \end{enumerate}
\end{lemma}

\begin{theorem}
    \label{ThmDilationOf3ConcaveWithoutExpansivity}
    Let $ H $ be a Hilbert space and let $ T\in \mathbf{B}(H) $ be $ 3 $-concave. Then $ T $ has a $ 3 $-isometric dilation $ W $ of the form \eqref{FormDilationMatrix} on the space $ H\oplus \ell^{2}(H^{\prime}) $, where $ H' \subset H $ is a closed subspace of $ H $.
\end{theorem}

\begin{proof}
    We provide only a sketch of the proof, because the construction of a dilation is very similar to the one given in the proof of Theorem \ref{ThmDilationOfMConcave}. Set $ \Delta = \beta_{2}(T) $. Define a sesquilinear form $ \psi\colon \mathcal{R}(\Delta^{\frac{1}{2}})\times \mathcal{R}(\Delta^{\frac{1}{2}}) \to \mathbb{C} $ by the formula:
    \begin{equation*}
        \psi(\Delta^{\frac{1}{2}}f,\Delta^{\frac{1}{2}}):= \langle T^{\ast}\beta_{3}(T)Tf,g\rangle, \quad f,g\in H.
    \end{equation*}
    Since
    \begin{equation*}
        T^{\ast}\beta_{3}(T)T = T^{\ast 2}\Delta T^{2} - T^{\ast}\Delta T
    \end{equation*}
    and, by 3-concavity,
    \begin{equation*}
        T^{\ast 2}\Delta T^{2} \le T^{\ast }\Delta T \le \Delta
    \end{equation*}
    we get that $ \psi $ is well-defined and continuous.  Extending $ \psi $ to a continuous sesquilinear form on $ \overline{\mathcal{R}(\Delta^{\frac{1}{2}})} $ and applying the Riesz theorem we obtain an operator $ A\in \mathbf{B}(\overline{\mathcal{R}(\Delta^{\frac{1}{2}})}) $ such that
    \begin{equation}
        \langle A\Delta^{\frac{1}{2}}f,\Delta^{\frac{1}{2}}g\rangle = \langle T^{\ast}\beta_{3}(T)Tf,g\rangle, \quad f,g\in H.
    \end{equation}
    It follows from the inequality $ \beta_{3}(T)\le 0 $ that $ A\le 0 $. Setting $ B = (I-A)^{\frac{1}{2}} $ and proceeding as in the proof of Theorem \ref{ThmDilationOfMConcave} we obtain a 3-isometric unilateral weighted shift with positive and invertible weights $ (S_{j})_{j\in \mathbb{N}_{1}} $ such that $ S_{1} = I $ and $ S_{2} = B $. Let $ W $ be as in \eqref{FormDilationMatrix} with $ U = \Delta^{\frac{1}{2}} $ and $ (S_{j})_{j\in \mathbb{N}_{1}} $ given above. By Lemma \ref{LemMIsometricityOfDilation}, it is enough to verify that \eqref{FormDilationMatrixMIsometric} holds to show that $ W $ is 3-isometric. For $ h\in H $ we have
    \begin{align*}
        &\langle \beta_{3}(T)h,h\rangle + \langle T^{\ast 2}\Delta T^{2}h,h\rangle -2\langle T^{\ast}\Delta Th,h\rangle + \langle S_{2}^{2}\Delta^{\frac{1}{2}}h,\Delta^{\frac{1}{2}}h\rangle\\
        &=\langle T^{\ast}\Delta Th,h\rangle - \langle \Delta h,h\rangle + \langle T^{\ast 2}\Delta T^{2}h,h\rangle\\ &-2\langle T^{\ast}\Delta Th,h\rangle + \langle S_{2}^{2}\Delta^{\frac{1}{2}}h,\Delta^{\frac{1}{2}}h\rangle\\
        &=\langle T^{\ast 2}\Delta T^{2}h,h\rangle - \langle T^{\ast}\Delta Th,h\rangle -\langle A\Delta^{\frac{1}{2}}h,\Delta^{\frac{1}{2}}h\rangle\\
        &=\langle T^{\ast}\beta_{3}(T)Th,h\rangle - \psi(\Delta^{\frac{1}{2}}h,\Delta^{\frac{1}{2}}h) = 0.
    \end{align*}
    Hence, \eqref{FormDilationMatrixMIsometric} holds. Therefore, $ W $ is a 3-isometric dilation of $ T $.
\end{proof}

\end{document}
